\begin{filecontents}[overwrite]{mir_fallback.bib}
% Fallback copies of canonical corpus entries, loaded after the shared refs.bib,
% which loads first and so wins every duplicate key (biber first-wins).
% Every key below is canonical in the corpus bibliography of September 19, 2026.
@article{Aethus,
  author      = {Benander, Ajax},
  title       = {Aethic Reasoning: A Comprehensive Solution to the Quantum Measurement Problem},
  year        = {2024},
  month       = {November},
  url         = {https://philpapers.org/rec/BENARA-5},
  eprint      = {https://philsci-archive.pitt.edu/id/eprint/25713},
  eprinttype  = {Alternative URL},
  note        = {Manuscript; the complete development of the framework, successive
                 deposits November 2024--2026. PhilArchive version history at
                 \url{https://philarchive.org/versions/BENARA-5}. This version: September 2026.}
}
@misc{AethusOptimalEarth,
  author       = {Benander, Ajax},
  title        = {The Optimal Earth: A Paradigmatic Re-Rendition of Natural History from Coupled Improbability},
  year         = {2026},
  howpublished = {Preprint},
  url          = {https://aethus.org},
  note         = {The natural-history extension; the conditional observer-mass
                  concentration theorem and the Optimal-Earth Conjecture.}
}
@article{Nexus,
  author      = {Benander, Ajax},
  title       = {Nexic Reasoning: Defining a Generalized Calculus Over Anthropic Parameters},
  year        = {2024},
  month       = {December},
  url         = {https://philpapers.org/rec/BENNRD},
  doi         = {10.31219/osf.io/jvkcp},
  note        = {Manuscript; the anthropic wing's master, successive deposits
                 December 2024--2026. Also at \url{https://osf.io/jvkcp}.
                 This version: September 2026.}
}

@misc{GoldenScience2022,
  author       = {Benander, Ajax},
  title        = {Golden Science},
  year         = {2022},
  month        = {November},
  howpublished = {Self-hosted manuscript},
  url          = {https://aethus.org/aethic-code/files/GoldenScience_Nov2022.pdf},
  note         = {Title page August 2022, the date of the first version of the source; this version created November 2022. Not deposited to a
                  preprint server; cited as a dated developmental record rather than
                  as a source of results.}
}
\end{filecontents}
\begin{filecontents}[overwrite]{mir_refs.bib}
% Policy: these locals cover what the shared refs.bib does not yet hold; refs.bib,
% loading first, wins any duplicate-key collision with its richer canonical
% entries.
% MIGRATE AT MERGE --- Golden-branch siblings, absent from the September 19, 2026
% canonical file. Add these to the shared refs.bib, whereupon (refs.bib loading
% first) the copies below become shadowed. The key SurvA, recorded under
% AethusGolden in that file, names GoldenSurvivorship, a distinct work.
@misc{GoldenSurvivorship,
  author       = {Benander, Ajax},
  title        = {Survivorship Conditioning for Hazard Rates Driven at the Log-Derivative Level},
  year         = {2026},
  howpublished = {Working draft},
  url          = {https://aethus.org},
  note         = {The Golden branch's root extraction; the two-condition viability
                  theorem with its piecewise driver-viability functional, the
                  hard-wall and soft-wall conditioned laws, the certification
                  ceiling, and the half-time theorem.}
}
@misc{GoldenTargeted,
  author       = {Benander, Ajax},
  title        = {Conditioning a Record on Its Own Future: Targeted Accordance and Survivorship-Conditioned Path Measures},
  year         = {2026},
  howpublished = {Working draft},
  url          = {https://aethus.org},
  note         = {The theorem extraction; future-targeted accordance under the
                  conditioned draw, the multiplicity guard, and the Bayes-factor
                  reading. Imports only GoldenSurvivorship.}
}
@misc{GoldenRouteWorld,
  author       = {Benander, Ajax},
  title        = {The Route World: A Golden-Reasoning Hypothesis of Biospheric Natural History},
  year         = {2026},
  howpublished = {Working draft},
  url          = {https://aethus.org},
  note         = {The terminal synthesis of the technical branch; the Golden Point,
                  the Departure Point, the Continuity Principle, and the
                  conditioned-record hypothesis with its prospective test.}
}
@misc{GoldenCivNature,
  author       = {Benander, Ajax},
  title        = {Is Civilization Anti-Nature? A Conditional Argument Against Capability Elimination, with Capability's Insufficiency and a Viability Constraint for Alignment},
  year         = {2026},
  howpublished = {Working draft},
  url          = {https://aethus.org},
  note         = {The application extraction; the capability-decay lemma, the
                  two-by-two viability picture, the winner's irony, and the
                  biospheric viability boundary for alignment.}
}

% External citations. Those not yet in the shared refs.bib should migrate there
% at the author's merge; duplicates resolve in refs.bib's favor.
% Author names corrected at the 2026-10-09 revision (Reuber, Victoria; Krug, Lena-Marie); VERIFY the full list against the publisher's page at submission.
@article{Burrow25,
  author  = {Pinkert, Stefan and Reuber, Victoria and Krug, Lena-Marie and Heidrich, Lea and Rehling, Finn and Brandl, Roland and Farwig, Nina},
  title   = {Burrowing facilitated the survival of mammals in harsh and fluctuating climates},
  journal = {Current Biology},
  volume  = {35},
  number  = {8},
  year    = {2025},
  pages   = {1779--1790.e3},
  doi     = {10.1016/j.cub.2025.02.064}
}
@article{During22,
  author  = {During, Melanie A. D. and others},
  title   = {The {M}esozoic terminated in boreal spring},
  journal = {Nature},
  volume  = {603},
  year    = {2022},
  pages   = {91--94}
}
@book{Goff23,
  author    = {Goff, Philip},
  title     = {Why? The Purpose of the Universe},
  publisher = {Oxford University Press},
  year      = {2023}
}
@book{Gross18,
  author    = {Gross, Alan G.},
  title     = {The Scientific Sublime: Popular Science Unravels the Mysteries of the Universe},
  publisher = {Oxford University Press},
  year      = {2018}
}
@article{Kaiho2017,
  author  = {Kaiho, Kunio and Oshima, Naga},
  title   = {Site of asteroid impact changed the history of life on {E}arth: the low probability of mass extinction},
  journal = {Scientific Reports},
  volume  = {7},
  year    = {2017},
  pages   = {14855}
}
@incollection{MetzSEP,
  author    = {Metz, Thaddeus},
  title     = {The meaning of life},
  booktitle = {The Stanford Encyclopedia of Philosophy},
  editor    = {Zalta, Edward N. and Nodelman, Uri},
  publisher = {Metaphysics Research Lab, Stanford University},
  year      = {2021},
  url       = {https://plato.stanford.edu/entries/life-meaning/},
  note      = {First published 2007; substantive revision 2021.}
}
@book{Mulgan15,
  author    = {Mulgan, Tim},
  title     = {Purpose in the Universe: The Moral and Metaphysical Case for Ananthropocentric Purposivism},
  publisher = {Oxford University Press},
  year      = {2015}
}
@article{Nature2025Sulfur,
  author  = {Rodiouchkina, Katerina and Goderis, Steven and Senel, Cem Berk and others},
  title   = {Reduced contribution of sulfur to the mass extinction associated with the {C}hicxulub impact event},
  journal = {Nature Communications},
  volume  = {16},
  year    = {2025},
  pages   = {620}
}
@book{Scheffler13,
  author    = {Scheffler, Samuel},
  title     = {Death and the Afterlife},
  publisher = {Oxford University Press},
  address   = {New York},
  year      = {2013}
}
@book{Weinberg77,
  author    = {Weinberg, Steven},
  title     = {The First Three Minutes},
  publisher = {Basic Books},
  year      = {1977}
}
@book{Wolf10,
  author    = {Wolf, Susan},
  title     = {Meaning in Life and Why It Matters},
  publisher = {Princeton University Press},
  year      = {2010}
}
@article{gardner1977mathematical,
  author  = {Gardner, Martin},
  title   = {Mathematical games: in which joining sets of points leads into diverse (and diverting) paths},
  journal = {Scientific American},
  volume  = {237},
  year    = {1977},
  pages   = {18--28}
}
@book{nietzsche1974gay,
  author    = {Nietzsche, Friedrich},
  title     = {The Gay Science},
  translator = {Kaufmann, Walter},
  publisher = {Vintage},
  year      = {1974},
  note      = {Original 1882.}
}
@book{Popper63,
  author    = {Popper, Karl R.},
  title     = {Conjectures and Refutations: The Growth of Scientific Knowledge},
  publisher = {Routledge and Kegan Paul},
  address   = {London},
  year      = {1963}
}
@article{Huemer07,
  author  = {Huemer, Michael},
  title   = {Compassionate phenomenal conservatism},
  journal = {Philosophy and Phenomenological Research},
  volume  = {74},
  number  = {1},
  year    = {2007},
  pages   = {30--55}
}
@article{Nagel71,
  author  = {Nagel, Thomas},
  title   = {The absurd},
  journal = {The Journal of Philosophy},
  volume  = {68},
  number  = {20},
  year    = {1971},
  pages   = {716--727}
}
@book{Moore1903,
  author    = {Moore, G. E.},
  title     = {Principia Ethica},
  publisher = {Cambridge University Press},
  address   = {Cambridge},
  year      = {1903},
  note      = {\S\S18--22, the principle of organic unities.}
}
@book{Kant1790,
  author     = {Kant, Immanuel},
  title      = {Critique of the Power of Judgment},
  translator = {Guyer, Paul and Matthews, Eric},
  publisher  = {Cambridge University Press},
  address    = {Cambridge},
  year       = {2000},
  note       = {Original 1790; the Analytic of the Sublime, \S\S23--29.}
}
@book{Goodenough98,
  author    = {Goodenough, Ursula},
  title     = {The Sacred Depths of Nature},
  publisher = {Oxford University Press},
  address   = {New York},
  year      = {1998}
}
@book{JasanoffKim15,
  editor    = {Jasanoff, Sheila and Kim, Sang-Hyun},
  title     = {Dreamscapes of Modernity: Sociotechnical Imaginaries and the Fabrication of Power},
  publisher = {University of Chicago Press},
  address   = {Chicago},
  year      = {2015}
}
@article{DArmsJacobson00,
  author  = {D'Arms, Justin and Jacobson, Daniel},
  title   = {The moralistic fallacy: on the `appropriateness' of emotions},
  journal = {Philosophy and Phenomenological Research},
  volume  = {61},
  number  = {1},
  year    = {2000},
  pages   = {65--90}
}
@book{Darwin1859,
  author    = {Darwin, Charles},
  title     = {On the Origin of Species by Means of Natural Selection},
  publisher = {John Murray},
  address   = {London},
  year      = {1859}
}
@incollection{Weber1919,
  author    = {Weber, Max},
  title     = {Science as a vocation},
  booktitle = {From Max Weber: Essays in Sociology},
  editor    = {Gerth, H. H. and Mills, C. Wright},
  publisher = {Oxford University Press},
  address   = {New York},
  year      = {1946},
  pages     = {129--156},
  note      = {Lecture delivered 1917; first published 1919.}
}
@book{Taylor07,
  author    = {Taylor, Charles},
  title     = {A Secular Age},
  publisher = {Harvard University Press},
  address   = {Cambridge, MA},
  year      = {2007}
}
@book{JosephsonStorm17,
  author    = {Josephson-Storm, Jason {\=A}.},
  title     = {The Myth of Disenchantment: Magic, Modernity, and the Birth of the Human Sciences},
  publisher = {University of Chicago Press},
  address   = {Chicago},
  year      = {2017}
}
@book{ListPettit11,
  author    = {List, Christian and Pettit, Philip},
  title     = {Group Agency: The Possibility, Design, and Status of Corporate Agents},
  publisher = {Oxford University Press},
  address   = {Oxford},
  year      = {2011}
}
@book{Christian04,
  author    = {Christian, David},
  title     = {Maps of Time: An Introduction to Big History},
  publisher = {University of California Press},
  address   = {Berkeley},
  year      = {2004}
}
@article{Swidler86,
  author  = {Swidler, Ann},
  title   = {Culture in action: symbols and strategies},
  journal = {American Sociological Review},
  volume  = {51},
  number  = {2},
  year    = {1986},
  pages   = {273--286}
}
\end{filecontents}
\documentclass[11pt]{article}
\usepackage[margin=1.1in]{geometry}
\usepackage{amsmath,amssymb,amsthm,mathtools}
\usepackage{mathrsfs}
\usepackage{booktabs,enumitem,url,xcolor}
\usepackage[colorlinks=true,linkcolor=blue!50!black,citecolor=blue!50!black,urlcolor=blue!50!black]{hyperref}
\usepackage[style=numeric, sorting=none]{biblatex}
% Load order: refs.bib first, so canonical entries win duplicate-key resolution (biber first-wins).
\addbibresource{refs.bib}
\addbibresource{mir_fallback.bib}
\addbibresource{mir_refs.bib}
\renewcommand*{\bibfont}{\small}
\newcommand{\E}{\mathbb{E}}
\newcommand{\PP}{\mathbb{P}}
\newcommand{\R}{\mathbb{R}}
\newcommand{\one}{\mathbf{1}}
\newcommand{\Gc}{\mathcal{G}}
\newcommand{\Ec}{\mathcal{E}}
\newcommand{\Mc}{\mathcal{M}}
\newcommand{\Sc}{\mathscr{S}}
\theoremstyle{plain}
\newtheorem{theorem}{Theorem}[section]
\newtheorem{proposition}[theorem]{Proposition}
\newtheorem{lemma}[theorem]{Lemma}
\newtheorem{corollary}[theorem]{Corollary}
\theoremstyle{definition}
\newtheorem{definition}[theorem]{Definition}
\newtheorem{premise}{Premise}
\newtheorem{hypothesis}[theorem]{Hypothesis}
\theoremstyle{remark}
\newtheorem{remark}[theorem]{Remark}
\title{Miracleism:\\ \large Exotic Structure and Meaning After Scientific Disenchantment}
\author{Ajax Benander\thanks{Email:
\href{mailto:akb9408@rit.edu}{akb9408@rit.edu}. This is a philosophical
paper that imports its technical results rather than re-proving them:
the Accordance Principle, the Rare-Earth inference, the High-Contrast
Earth Hypothesis, and Optimal Earth from the Nexic companions
\cite{Nexus,AethusOptimalEarth}; the survivorship model and its ceiling from
\cite{GoldenSurvivorship}; future-targeted accordance from \cite{GoldenTargeted}; and the
Golden Point and the Route World hypothesis from \cite{GoldenRouteWorld}. Its
own contribution is the interpretive argument, and every claim it makes
is graded. Drafting collaboration: Anthropic Claude (Opus~4.x,
Fable~5), under the author's direction; the first-person sites remain
the author's to write.}}
\date{October 2026 --- Working Draft, rev.~8.2}
\begin{document}
\maketitle

\begin{abstract}
Scientific inquiry has undermined some traditional cosmological warrants for
human significance, and self-authored purpose cannot, by itself, reproduce
the discovered cosmological warrant at issue, because it preserves the feature
the loss exposed: the warrant's source remains ourselves. This paper asks whether further inquiry could
disclose an independently existing structure that supplies a different
warrant, and argues, conditionally, that it can. If the
realized natural history occupies an extraordinarily narrow and
target-accordant region of possibility space, and if that region is a coupled
structure of which the observer is a physical continuation, then there is an
independently discovered object to which existential significance can answer,
and the significance is participatory: the discoverer is part of what is
discovered, at a point where its continuation can still be preserved or
forfeited. The paper's own contribution is the argument from that relation to
two defeasible verdicts, that awe is fitting toward the object and that its
continuation gives its inheritors a reason for care, defended against a reader
who grants the descriptive claims and disputes the reasons. The end-Cretaceous
bottleneck is examined as an inheritance that is at once valuable and tragic,
and the Golden Point as the condition under which inherited significance
reaches living fruition without the future creating or erasing its source.
Every empirical premise is imported at its stated grade from the companion
papers; the bridges are philosophical premises and are marked as such.
\end{abstract}

\section{The question disenchantment leaves}\label{sec:valley}

Two announcements frame the strand of modern thought this paper begins from.
Nietzsche's, in 1882, was that the source from which meaning had been received
could no longer be believed in \cite{nietzsche1974gay}; Weinberg's, in 1977,
from inside the physics that had replaced it, was that the more the universe
seems comprehensible the more it seems pointless \cite{Weinberg77}. Between
them lies a familiar genealogy: explanation displaced the special creation,
the central Earth, and the designed body, and each displacement was correct
because the explanation was better than the object it replaced. Nietzsche
reported from inside what it is like when a whole structure that had grounded
values fails at once; Weinberg's sentence extrapolates the slope he could see.
Neither is corrected here. On the evidence available to them, the downward
inference was reasonable.

The genealogy is a schematic account of one strand, offered as motivation, and
it is narrower than three things it is easily mistaken for. It is not a
history of secularization: Weber's disenchantment concerned the retreat of
magic and of ultimate meaning from a rationalized world, the loss of a
cosmological warrant being one thread among several \cite{Weber1919};
Taylor's secular age is a history of the conditions of belief, in which belief
and unbelief became options within a shared immanent frame, not a story of
subtraction \cite{Taylor07}; and Josephson-Storm argues that the
disenchantment narrative is itself partly a myth, enchantment having persisted
alongside and inside modern science \cite{JosephsonStorm17}. Nor does the scientific
displacement of particular theological explanations establish nihilism, a
philosophical conclusion with its own arguments; Weinberg's sentence reports a
mood as much as an inference. And science has
not previously produced only disenchantment. Darwin closed the \emph{Origin}
by finding grandeur in the view of life his explanation afforded
\cite{Darwin1859}, and the scientific sublime has a literature of its own
\cite{Gross18}. That predecessor is a resource for this paper rather than an
inconvenience: naturalistic wonder is older than Miracleism, and nothing below
claims otherwise.

What the strand does establish is narrower and sufficient. The objects
previously offered as meaning's external warrant were removed by explanation;
the removal was a gain in understanding; and a warrant of the removed kind
cannot be restored by lowering the standards that removed it. Returning to the
certainty of the premodern peak would be a return by discarding the knowledge
responsible for the descent, which is why the paper sides with the settlement
against every restoration; the dragon that supplied meaning by coupling fear
with control stays dead. The live question is whether the same engine, run
further, finds anything else.

The question needs to be asked precisely, because several questions travel
together under the word \emph{meaning} and their answers do not arrive
together. One is descriptive: what kind of natural history produced us. One
concerns fitting emotion: whether that history is an appropriate object of
awe. One concerns existential significance: whether our participation in it
makes our lives significant. One is practical: whether we have reasons to
preserve its continuation. One is social: whether a shared understanding of it
could change how a species deliberates. An object can warrant awe without
warranting preservation; a person can participate in something significant
without special responsibility for its future; a true account can be
motivationally inert. This paper's thesis concerns the third and fourth
questions, and it concerns a particular form of significance, which may be
called \emph{participatory}: the significance available to an agent who
discovers an independently existing natural history, recognizes itself as a
continuation of that history, and understands its present agency as bearing on
the history's unfinished future. Three things that can come apart are placed
now. Significance attaches through participation itself: being a
continuation state of the history is a fact about the agent that holds
whether or not it is recognized, and it is the ground of everything that
follows. Recognition of that participation is what makes awe fitting
\emph{for} the agent and gives the agent access to the reasons the relation
supplies. Engagement with the continuation is where the reason for care has
its object and where responsibility attaches. Participatory significance is
therefore not a third verdict beside awe and care; it is what both are about,
the fact of participation supplying the object, its recognition the
fittingness, and its continuation the practical reason.

At each transition below the premise doing the work is named, and
Section~\ref{sec:ladder} collects the premises with their grades.

One commitment is announced now. The companions supply two descriptive claims of different
strength. The weaker is that the record is rare and target-enriched: the
realized history is improbable under the baseline and enormously overweighted
under conditioning on the biosphere's survival. The stronger, the Route World
hypothesis, is that the law governing the realized historical ensemble is the
conditioned law itself, so that target-consistent structure is generated
rather than retrospectively selected. Much of the interpretation below is
available under the weaker claim alone; what the stronger adds is stated where
it enters (Sections~\ref{sec:kpg} and~\ref{sec:golden}), so that the
contribution of the stronger ontology is visible without the argument
depending on it everywhere.

\section{The object: rarity}\label{sec:rarity}

The evidence begins in probability, and probability is the doorway and not
the final ontology. The Nexic companion's Accordance Principle is the
constraint Bayes' theorem returns from a Copernican draw together with the
existence of a conditioning target: a realized attribute's prior improbability
is compensated by its relevance to the target, with $\kappa$-fold violations
confined to probability $1/(1+\kappa)$ \cite{Nexus,GoldenTargeted}. Run on the
record, the principle yields the Rare-Earth inference under stated antecedents
and, with the High-Contrast Earth Hypothesis, the Optimal-Earth program: the
realized natural history is not one ordinary member of a large equivalence
class of living worlds but sits in an extraordinarily narrow region of the
space of possible histories, and in a region the target picks out
\cite{AethusOptimalEarth}.

\begin{definition}[The two magnitudes]\label{def:magnitudes}
Let $A^\ast$ be an independently specified configuration of the record,
$C$ the common baseline of observers existing now and the biosphere
having survived to the present, and $G_T=\{T_{\mathrm{dep}}>T\}$ the
horizon event that the biosphere survives to horizon $T$, the target
being read through the horizon family because in the critical regime
the event of indefinite survival is null and cannot be conditioned on
directly \cite{GoldenSurvivorship,GoldenTargeted}. The \emph{surprisal} of $A^\ast$ is
\[
I(A^\ast)=-\log\PP(A^\ast\mid C),
\]
and its \emph{selection amplification} at horizon $T$ is
\[
\Mc_T(A^\ast;\Gc)=\log\frac{\PP(A^\ast\mid G_T,C)}{\PP(A^\ast\mid C)}
=\log\frac{\PP(G_T\mid A^\ast,C)}{\PP(G_T\mid C)},
\]
the second equality by Bayes at every finite $T$; where the companions
have established the limiting conditioned law, the amplification under
the target is $\Mc(A^\ast;\Gc)=\lim_{T\to\infty}\Mc_T(A^\ast;\Gc)$, and
every expression below written with $\Gc$ is shorthand for that limit
where it exists and for the finite-horizon form otherwise. A
configuration is a candidate Miracle when both magnitudes are large:
rare under the baseline, and enormously enriched by the target the
history appears to serve.
\end{definition}

The distinction defeats the cheapest objection, that every sufficiently
detailed history is improbable. The shuffled deck has $I$ of about $226$ bits
and $\Mc$ of zero for any target independent of the deck's ordering, since
nothing about such a target enriches one ordering over another, and a target
invented to prefer a particular permutation fails the independent-specification
guard below. A Miracle in the sense this paper needs is a rare event that the
target amplifies, and the amplification is a statement about the record and
the target, not about the fineness of a description.

Four guards keep the magnitudes from being manufactured, and they are
requirements on any configuration the paper calls a Miracle. The attribute
must be \emph{specified independently} of its rarity, before or without
reference to the record's detail, so that no retrospective microdescription
inflates $I$; a configuration built from every molecule and every coin toss
has probability near zero and signifies nothing. It must be
\emph{counterfactually relevant} to the target, a requirement separate from
enrichment: $\Mc_T$ measures association under the specified law, the
overweighting of $A^\ast$ in the conditioned ensemble, and does not say what
would have happened to the target had $A^\ast$ been otherwise. The
counterfactual test is the Route World companion's route-deletion audit, read
under the framework's native settling operation, in which an attribute is
settled to its alternative and the target's probability is recomputed on the
settled state \cite{GoldenRouteWorld,Aethus}. The two coincide only where
settling $A^\ast$ leaves every other entry's standing unchanged; in general
settling detaches an attribute from its grounds while conditioning does not,
and no reduction of the one to the other is assumed. Where only enrichment is
available, the configuration is a candidate and not yet a Miracle. Its rarity must be \emph{robust across a defensible model class}, so
that the reported magnitude is a conservative lower bound,
\[
I_{\min}(A^\ast)=\inf_{\theta\in\Theta}\bigl[-\log\PP_\theta(A^\ast\mid C)\bigr],
\qquad
\Mc_{\min}(A^\ast;\Gc)=\inf_{\theta\in\Theta}\log\frac{\PP_\theta(A^\ast\mid\Gc,C)}{\PP_\theta(A^\ast\mid C)},
\]
taken over models deliberately chosen to make the event ordinary. And
wherever possible the configuration should carry \emph{prospective}
confirmation, a prediction scored on record not yet read, which the Route
World companion's protocol supplies. The paper claims no particular figure
such as $10^{-1000}$, and it does not re-derive the companions' bounds. Its
claim is conditional on them, and the results consumed are these, each at
the companions' deposited version of September 2026. The exceedance pricing
is the theorem companion's Theorem \emph{Future-targeted accordance} with
its multiplicity guard, in the section of that name \cite{GoldenTargeted};
its assumptions being a draw under the target-conditioned measure and, for
survival targets, the survivorship model's horizon family. The route-deletion requirement is the
Route World companion's Proposition \emph{Significance from route-deletion}
and the remark on the counterfactuals it excludes, in the subsection on
flukes \cite{GoldenRouteWorld}, stated there within the model and graded
structural rather than proved. The Rare-Earth inference is the Nexic
companion's Accordance Principle with its Substitution Property and
Common-Ground Schema, which bound the abundance of any articulable
observer-class by inclusion; the quantitative figure is that companion's
Dinosaurian-extrapolation corollary in the subsection \emph{Followup Nexic
Extrapolations}, an upper bound of order $10^{-27}$ per world, of order one
observer-bearing world per observable volume, under the stated assumption
that rare target-relevant attributes are not anomalously concentrated in the
interval analyzed, which the companion itself grades below the Accordance
Principle \cite{Nexus}.
The robustness requirement displayed above, the infimum over a model class
chosen to make the event ordinary, is this paper's requirement on any
candidate; no companion result is claimed to have discharged it, and where
the companions' assumptions are contested the claim is contested with them.

One further result is stated because it licenses a phrase that is easily
over-read; it concerns records with a certain structure and asserts of no
record that it has it.

\begin{proposition}[Unbounded amplification along a joint limit; conditional]\label{prop:nested}
If the record admits an independently specified nested family of
canonical partitions $A_1\supset A_2\supset A_3\supset\cdots$ for which $\PP(A_n \mid C) \to 0$ while, along a joint limit $T_n \to \infty$ with $\PP(G_{T_n} \mid C) \to 0$, $\PP(A_n \mid G_{T_n}, C)$ stays bounded away from zero, then $\Mc_{T_n}(A_n;\Gc) \to \infty$ as $n \to \infty$; at any fixed finite horizon $T$ with $\PP(G_T \mid C) > 0$ the amplification is bounded by $-\log \PP(G_T \mid C)$, a ceiling independent of descriptive refinement, so the absence of a ceiling is a statement about the joint limit and the order in which it is taken, never about a fixed horizon:
the evidential contrast the method extracts has no ceiling set by the
method, only by the record.
\end{proposition}

\noindent The proposition's content is unbounded evidential contrast under
canonical refinement, and nothing more: no finite record carries infinite
evidence. Whether the record of natural history admits such a family is an
empirical question the Nexic program addresses and this paper does not
settle; the proposition licenses no numerical lower bound for the actual
history, and the canonical-partition requirement keeps the family from being
manufactured by redescription.

Rarity does two jobs. The first is epistemic: it is what makes an observer
stop and say that something requiring explanation is happening; without it,
coupling looks like ordinary causal organization. The second is cosmic, and
it is stated in two steps. The Rare-Earth inference, under the Nexic
companion's abundance bridge named above, yields a low density of comparable
observer-bearing biospheres. Low density is not yet a small number: the step
from density to count requires a domain, the observable volume under the
cosmic weighting, and the quantitative bound named above, imported without
re-derivation \cite{Nexus}. ``One of extraordinarily few'' is therefore a
conditional claim, relative to that domain and to that bound. If the bound
holds, the reader's indexical situation changes, from one observer-bearing
world among multitudes to one of few within the volume that could ever be
surveyed. The stronger image, that the expanse is
sterile of life altogether, needs an additional bridge from observer rarity
to biotic rarity, since simple or even complex non-observer life could be
common while observer-bearing histories remain rare; that bridge is a
separate empirical conjecture. The conservative version is what the paper
relies on; the ambitious version is marked where it is used.

What the magnitudes deliver, before any interpretation, is a scale conveyed
by the only analogy that fits: Graham's number, rigorously defined by a
compact recursive notation, can be established to lie beyond ordinary
numerical representation while no proportionate intuitive representation of
it can be formed \cite{gardner1977mathematical}. The proposition that this
route occupied an extraordinarily tiny fraction of baseline possibility space
can be fully understood as a fact and remain unrenderable.

\section{The object: coupling and exotic structure}\label{sec:exotic}

A number, however small, is inert. $10^{-1000}$ by itself does not feel like
anything, and if Miracleism were the sublimity of improbability it would be
answered by the shuffled deck and by the observation that awe at a small
number is a category error. The thing encountered is upstream of the number.

The Nexic companion's hierarchy says where: a flat probability is a lossy
projection of a hazard and improbability structure, itself a lossy projection
of an Aethus-valued structure, the record as a set of stated attributes with
their dependencies \cite{Nexus,Aethus}, so that reducing a Miracle to
$\PP(A)\ll1$ discards the feature the paper finds significant. The Accordance
constraint does not terminate where it is applied: one extraordinary
requirement implies another, which constrains a third, which explains why a
fourth apparently bizarre feature had to occupy its particular place, and the
closure of that propagation is a single object rather than an inventory of
independent surprises,
\[
\Ec=\operatorname{Closure}(A_{\mathrm{root}};\ \text{Accordance, coupling}),
\qquad
\Ec\neq\{A_1,A_2,\dots,A_n\}.
\]
The object is given witnesses so that coupling and depth are inputs and not
adjectives; the formal apparatus and its audits are in
Appendix~\ref{app:witnesses}. In outline: the witnesses are native to the
framework's settling semantics, the relation that matters being what happens
to an attribute's target effect when other attributes are jointly settled to
their alternatives, and generic statistical dependence is not the measure.
The dependency object is a weighted interaction hypergraph whose hyperedges
record higher-order target interactions that no pairwise test can see; the
exotic structure generated by a root is that hypergraph together with its
forcing relations on the closure the root generates, a structured closure
and not the bare inventory of its members, which is what the inequality
above marks. Its witnesses are a coupling mass, a depth, and the rarity of
the realized conjunction. \emph{Non-reductive} means nonseparability of the
specified kind and nothing more; depth is the iteration depth of a stated
presentation and not an intrinsic height, and the argument uses of it only
that the closure is not reached at the root; and the witnesses are defined
on description classes modulo refinements carrying no target information,
with the invariance that is constructed and the invariance that is required
kept apart. The definition is an interpretive formalization of conjectural
upstream structure, graded \textbf{[I]} on the companions' \textbf{[C]}.

A mechanism with no forcing relations and no target interaction at any
admissible order would be reductive: its interaction hypergraph empty, its
depth zero, its identity the list of its members. Exotic structure is not
that. Its identity resides partly in the relations, so that perturbing a major
member reconfigures the rest and destroying the coupling destroys some of what
the thing is rather than only the bookkeeping about it; the Nexic companion
calls the optimal Aethus a knot, its parts standing where they do because the
others do. Rarity supplies the contrast with the baseline and coupling the
integration that makes the closure one object rather than a list, and nothing
in either is yet a claim about value, which is the business of
Sections~\ref{sec:construction} and~\ref{sec:sublime}. Rarity is exotic
structure's scalar shadow, not its definition: the shuffled deck is low
probability with low coupling, no depth, and no inherited relation to the
observer, and no scalar cutoff on probability makes anything exotic.

\section{Provenance and value}\label{sec:construction}

One influential secular response to the loss, the constructionist settlement,
is that whatever meaning there is, we make. It is not the only position
philosophy has taken: naturalistic objectivist accounts hold that
mind-independent features of the world can confer meaning, and hybrid accounts
such as Wolf's hold that objective attractiveness must meet subjective
attraction \cite{Wolf10,MetzSEP}. Miracleism claims no priority for the bare
possibility of naturalistic objective meaning. Its proposed contribution is
the object alleged to supply it and the observer's relation to that object.
The settlement is granted its premise entirely; what is not granted is that
construction can reverse the descent, and the reason is a matter of
provenance.

Consider a person who begins with no externally licensed meaning and some
structure of concerns, and who perturbs it at will: a career, a cause, a
community, a moral theory, humanity, art, the progress of science. The
projects can be made as noble as one likes, and every such transformation
preserves one property, the one the descent removed. Ask of the candidate
warrant behind a meaning structure, rather than of the significance itself,
where it came from: whether it was authored by the agent or the human
community, or encountered in the world and not put there. Self-authored
transformation of self-authored warrant yields self-authored warrant, and the
statement is close to definitional, since the transformations were defined as
self-authored; that is its strength and its limit, and it is not a proof that
constructed meaning is not meaning. It
is the narrower observation that self-authorship cannot reproduce the
particular external provenance that disenchantment removed, which was not a
purpose of any particular size but a purpose of the other provenance class:
something the person did not put into the world and encountered as larger than
the scale on which they ordinarily manufacture goals. What could change the
provenance class is not another project. It is evidence that changes what the
world itself warrants.

Provenance cannot carry the evaluative burden by itself, however, and three
questions that the word \emph{independent} runs together must be separated.
Whether the object exists independently of our wish for significance. Whether
our evidence about it is independent of that wish. Whether its existence
supplies reasons to value it. The first two establish that Miracleism is
responsive to reality rather than constitutive of its object, and the
companions' protocols are addressed to them. The third is the philosophical
claim, and two ordinary cases show that it does not follow from the first
two. A human-created work of art can possess merits its maker cannot determine
by wishing; authorship does not make value arbitrary. A naturally occurring
destructive process, a pandemic strain or an advancing ice sheet, is entirely
independent of human authorship and deserves no continuation; external
provenance does not make a thing valuable. The provenance distinction
identifies a difference in source. Which sources confer value is a separate
question, and Section~\ref{sec:sublime} answers it for one object by argument
rather than by provenance.

The same discrimination applies to the taxonomy of how existential authority
reaches a person. \emph{Authored}: I supply the significance.
\emph{Received}: another alleged authority tells me the significance.
\emph{Discovered}: the world supplies evidence from which I find it. The
premodern religious source is predominantly received, the constructionist
reconstruction predominantly authored, and the third cell's defining property
is that the source of warrant is an independently encountered natural fact,
so that valuation is responsive to the object. But the cells are not
watertight: chosen projects and relationships encounter standards and reasons
their participants did not invent. What the discovered cell adds is the independence of the object
itself: a natural history that would have the structure it has whether or not
anyone found it worth engaging.

Wolf is the right interlocutor here, because her account already joins
subjective engagement to objective worth and so offers a far stronger
alternative than arbitrary self-authorship: meaning arises, on her view, from
active engagement in projects of worth, the worth not conferred by the
engagement \cite{Wolf10}. What Miracleism's object adds to that relationship
has three parts, of which the paper defends the first two and holds the third
at the stronger grade. It proposes
a special relation to an existing kind of value: the goods of living history,
experience, understanding, relationship, and the biosphere's living forms, are
recognizable objects of worth on many accounts of objective worth, and
the discovery relates the reader to them as their inheritor and continuation
rather than as their admirer. It proposes a new reason for engagement: the
continuation of those goods is, on the companions' claims, irreplaceable and
still open, which changes what engagement with them is for. And it proposes,
at the stronger grade, a new bearer of worth, the integrated history itself,
whose value is not exhausted by the goods it carries; Section~\ref{sec:sublime}
says what that claim requires. On Wolf's scheme the object supplies the worth
and the discovery supplies a route to the attraction, one that is responsive
rather than constitutive.

Transmission is not authorship. A reader can receive testimony about a
discovered warrant, as nearly everyone receives most of what they know of
natural history, and communities can author educational materials and
practices responsive to it, without the underlying object becoming a product
of the demand it answers; the theory of evolution reached most people through
testimony and is not thereby authored by them. The paper's claim concerns the
source. Acquisition and expression are the subject of
Section~\ref{sec:selfaware}.

\section{Awe and care}\label{sec:sublime}

What the discovery, should it be secured, does to a person is answered by
decomposition rather than by definition, because the phenomenon is partly
ostensive. The decomposition is an inventory and not a response function: no
map from its components to any feeling is claimed.

\subsection*{The Nexic sublime}

The encounter with exotic structure has five components, four from the Nexic
companion and one from the Golden companion.

\begin{center}\small
\begin{tabular}{@{}llp{0.56\linewidth}@{}}
\toprule
Component & Witness & What it contributes to the judgment\\
\midrule
Rarity $R$ & $I_{\min}$; $R(\Ec)$ & scarcity, and under the Rare-Earth bridge the scarcity of comparable observer-bearing worlds\\
Coupling $K$ & $K(\Ec)$ & integration: the object is a knot and not a list\\
Depth $D$ & $D(\Ec)$ & recursive articulation: resolving one node can expose further coupled prerequisites beneath it\\
Incorporation $X$ & descent & the observer is a continuation state of the structure observed\\
Futurity $F$ & $G_T$ & the continuation is unfinished and can still be reached or forfeited\\
\bottomrule
\end{tabular}
\end{center}

\noindent Scale and depth are the stuff of the mathematical sublime, and
integration is what distinguishes an organism from a heap; incorporation and
futurity are what distinguish the Nexic sublime, and they are best seen
against Kant. On Kant's account the sublime is occasioned by the object's
defeat of the imagination and consists in the mind's discovery of a capacity, reason's
idea of totality or the moral vocation, that the sensible magnitude cannot
touch, so that the subject is elevated above nature in the very experience of
being overwhelmed by it \cite{Kant1790}. The Nexic sublime shares the first
movement, the magnitude that cannot be rendered, and reverses the second. What
the encounter discloses is not a capacity that transcends the object but a
constitution that belongs to it: the present biosphere is one of the
structure's continuation states, and so is the reader,
\[
A_{\mathrm{abiogenesis}}\preceq A_{\mathrm{Cambrian}}\preceq
A_{\mathrm{end\text{-}Cretaceous}}\preceq\cdots\preceq A_{\mathrm{now}}.
\]
The thing whose extraordinary structure is being discovered includes the
discoverer. The disclosure reconciles rather than elevates; it is the
discovery of dependence, and the observer's own vulnerability enters the
fitting response, because the vulnerability is the structure's, shared by
every continuation state that has so far survived. And because the object is
unfinished, contemplation is not sealed off from action: to apprehend the
structure adequately is to apprehend oneself at a point from which its
continuation can be affected, which no mountain or storm asks of its
spectator. This is the content of the phrase that stands at the literary
center of the paper, \emph{too vast to be conveyed and too personal to be
named}: vast warranted by rarity, coupling, and cosmic scarcity; personal
warranted by descent, the reader physically downstream of exactly the route
the exhibit of Section~\ref{sec:kpg} describes; and the failure of language
to convey their conjunction part of the phenomenon rather than a defect of
the argument.

\subsection*{Fitting awe}

Miracleism makes a stronger claim than that awe is available: conditional on
the empirical ladder holding, and on the object being apprehended at the grain
the preceding sections specify, awe is the fitting response to it.

\begin{premise}[Fitting-awe bridge; \textbf{[$\Phi$]}]\label{prem:awe}
Awe is defeasibly fitting toward an object when adequate apprehension
reveals a conjunction of cognitive vastness or depth, non-reductive
structured unity, and an indexically consequential relation between the
observer and stakes extending beyond the observer's ordinary local
horizon.
\end{premise}

\noindent Fittingness is a claim about the accuracy of an emotion's
presentation of its object, and it should be distinguished from two things
with which the word \emph{appropriate} is often conflated: from the moral or
prudential appropriateness of having the emotion, and from the probability
that a person will feel it \cite{DArmsJacobson00}. Awe presents its object as
vast or deep beyond the subject's ready grasp, as unified rather than
scattered, and as bearing on the subject in a way that calls for accommodation
rather than mere notice. If the ladder holds, the object has those properties:
rarity and depth supply the vastness, coupling the unity, incorporation and
futurity the bearing. So awe is fitting in the sense that its evaluative
presentation of the object is accurate. Nothing in that claim compels a
nervous system, and affect supplies no evidence for the framework; a reader
who finds in themselves curiosity, grief, or ambivalence is not thereby
mistaken, since fittingness is non-exclusive. What the bridge denies is narrower: a response that registers only
the scalar probability has not yet taken in the object described, since the
object is not the probability. Felt awe is no test of comprehension;
comprehension is a test of what one's awe, if any, is about.

Four attitudes should then be kept apart, because the argument's most tempting
error is to let one carry the others. Awe is fitting to scale, unity, and
bearing, and a threatening object can be awe-inspiring. Admiration is fitting
to achievement or excellence, and sits uneasily here, since no agent achieved
the structure; what can be admired without an agent is the form of a solution,
the way the end-Cretaceous configuration threads a joint constraint, and that
admiration is compatible with horror at what the solution cost. Affirmation is
a verdict that the whole is good, or that one is glad it occurred, and it is
not entailed by awe or admiration; Section~\ref{sec:kpg} argues that it must
be issued, if at all, without retrospective approval of every cause. Care is a
practical attitude with reasons of its own, and it is the subject of the
second bridge. Awe supplies none of the practical content; a person can be
awed by the structure and conclude nothing about what to do. The transition
from awe to care is therefore not a transition at all. It is a second
premise.

\subsection*{The care bridge, and the opponent who grants the science}

\begin{premise}[Care bridge; \textbf{[$\Phi$]}]\label{prem:care}
If an agent adequately apprehends itself as a continuation of a
non-authored, extraordinarily scarce, internally structured history
carrying goods worth sustaining, whose live continuation can still be
preserved or forfeited, that relation supplies a defeasible reason to care
about the continuation. It does not by itself supply a categorical duty,
rank competing reasons, or imply that every continuation is good overall.
\end{premise}

\noindent The premise is the value-theoretic step the probability theory
cannot perform, and it needs a defense against a reader stronger than the
shuffled-deck objector. Grant this reader everything descriptive: the history
is exceptionally rare, its structure internally coupled, we are continuations
of it, its future remains consequential, and even the stronger claim, that
the governing law is conditioned on the continuation, is true. The reader
asks why any of this makes the history valuable rather than merely
remarkable, and why being its continuation gives them a reason to preserve
it. They are not protecting nihilism from evidence; they dispute the account
of reasons, and Nagel supplies their sharpest form: being part of something
larger, longer, or more encompassing does not by itself answer the demand for
justification, since the larger thing's mattering can be questioned exactly as
one's own was \cite{Nagel71}. Magnitude, duration, and enclosure within a
vast project are not reasons. The care bridge must say what the particular
relation between the observer and the history contributes beyond them.

The answer assigns different features different work, and the thesis is the
conjunction. Begin with what is valuable. Nothing in the account asks anyone
to value improbability. The history whose continuation is at issue is the
realized history of goods that many accounts of objective worth already
recognize and that the reader addressed here grants: living experience,
understanding, relationship, creativity, and
the living forms of the biosphere in which those goods have so far been
carried. Those goods generate reasons for care wherever they occur, and the
first thing the discovery does is show that they are carried, here, by one
integrated and vulnerable continuation rather than by a scattering of
replaceable instances. Scarcity then does not create value; it changes the
reason's object. In a world where living worlds were abundant, the reason to
care for the goods would attach to the goods wherever found, and the loss of
one carrier would be compensable by others; where the carriers are, within the
surveyable volume, extraordinarily few, the reason attaches to this
continuation, which has no substitute. That is the difference between a good
and an irreplaceable good, and it is a difference in reasons, not in feeling.
Incorporation adds a relational reason of a kind ethics already recognizes.
Reasons arising from relationships one did not choose, to a family, a
language, a received body of knowledge, are not consented to and are not
thereby void; one can be bound to what one is made of. The observer is not a
spectator who has learned that a rare thing exists somewhere. The observer is
a continuation state of it, constituted by its passage through the
bottlenecks, and the reason to care for its continuation is of the same kind
as the reason a person has to care for what has made them what they are, with
the difference that here the relation is literal and physical. Nagel's
challenge is answered at this point and not before: what the relation
contributes is not size but constitution. Futurity adds the prospective
element. The continuation can still be preserved or forfeited, and agents now
have leverage over which; responsibility attaches where relationship and
situated agency meet, and its extent, which Section~\ref{sec:selfaware}
returns to, depends on what an agent can reasonably do.

Two routes to the bridge have now been distinguished, and the choice between
them is stated. One holds that the unusual structure itself has evaluative
importance; the other begins with recognizable goods and argues that the
discovery reveals their integrated, vulnerable, and irreplaceable
continuation, rarity intensifying an existing reason rather than creating
value from nothing. The argument above runs the second route, which suffices
for every conditional conclusion the paper draws about awe, care, and
responsibility. The author's thesis is the stronger one, and it should be stated
precisely, since two claims travel under that description. The weaker of the
two is that the history's organization makes the goods it carries a
particular, nonfungible object of concern; that claim follows from the
identity condition of Section~\ref{sec:exotic} and is already part of the
goods route. The stronger, which is the thesis, is that the integrated
history has value beyond the independently valuable goods it carries. Identity
alone does not establish it: a harmful institution can possess a distinctive,
integrated, irreplaceable identity, and its loss would remove a particular
thing, without that thing meriting valuation. The positive evaluative premise
is therefore this. The organization of this history is an organization
\emph{of goods}: the goods it carries are its products and its mode of
continuation, their identities include their places in it, and the relations
that integrate them are the relations by which they came to be and are
sustained. An integration of goods of that kind is a bearer of value in its
own right, in the way a life is valued as a whole beyond the sum of its good
moments, or a work beyond the sum of its passages, where destroying the unity
while preserving the parts' value sum is itself a loss. The premise is an
organic-unity premise restricted to unities of goods, and it is stated for
mixed histories, since no actual history is a unity of goods alone. A harmful
institution can produce and sustain some genuine goods, and the natural
history at issue contains suffering and destruction beside what it carried
forward, so the premise does not sort organizations into the wholly
beneficial and the wholly harmful. It says that the relations by which goods
are produced and sustained contribute positive value in proportion to the
goods they organize, and that the contribution neither makes the whole good
overall nor justifies its costs. A mixed history can bear organizational
value and remain tragic, the verdict Section~\ref{sec:kpg} issues; an
institution that organizes few goods and much harm bears little of it and is
defeated on balance, the hostile-inheritance case below. The premise adds a
value to be weighed, not a verdict that preempts weighing, which is what the
care bridge's defeasibility already provides for.

The premise has a name, and the comparison fixes what it adds. Moore's
principle of organic unities holds that the value of a whole need not equal
the sum of the values of its parts \cite{Moore1903}, and Moore separated
that claim from causal interdependence among the parts: a whole can be
non-additive in value whether or not its parts depend on one another, and
parts can depend on one another without the whole bearing value beyond
their sum. That separation is exactly the gap the stronger thesis must
cross. Section~\ref{sec:exotic} and Appendix~\ref{app:witnesses} establish,
conditionally, the target-level interdependence of the history's members;
the organic-unity premise is the further evaluative commitment that this
interdependence, being an organization of goods by the relations that
produced and sustain them, is one of the wholes whose value exceeds the sum.
Moore's principle licenses the possibility of such wholes and does not
establish that any coupled history is one; the premise asserts it for this
object, at the grade of a philosophical premise.

What the stronger claim adds to the goods route is then determinate, and the
word \emph{uncompensated} is given its scope. On the goods route, comparable
goods elsewhere compensate the loss of this carrier, and what scarcity adds
is urgency. On the stronger claim, relocating comparable goods does not
preserve this organization of them, so the loss of the whole is not made
good by any relocation of its goods; that is the claim made. The broader
claim, that the loss could not be compensated in any evaluative respect, is
not made: a successor whole might bear organizational value of its own, and
whether it would outweigh the loss is a comparison of wholes the premise
does not settle. Fruition accordingly concerns the continuation of this
organization rather than the survival of goods of its kind.

Four things should accordingly be kept apart, and the abundant-world case
below is read through them: replacement of comparable goods, which an
abundant world permits; replacement of the particular history, which no
world permits; compensation for its loss, which the goods route allows in
part and the stronger claim denies for the organization itself; and the additional importance scarcity
gives to preserving it, which abundance reduces. Abundance could reduce the
scarcity-based reasons without making this history replaceable or its loss
made good by relocating its goods. The two routes are therefore layered rather than
rivals: the goods route says what is valuable; the structural route says why
this organization of those goods, as this organization, is the object and why
its loss is a loss of its own. A reader who rejects the organic-unity premise
keeps the paper's conditional conclusions; a reader who accepts it has the
author's thesis.

\subsection*{Three cases}

Three cases isolate what each feature contributes, and none needs the full
Golden hypothesis. \emph{An abundant world.} Suppose life and worthwhile
experience are common across the surveyable volume. The goods route and the
care bridge's reason survive, and comparable goods elsewhere could replace
the goods this carrier holds, so the scarcity-based reasons weaken and the
reason detaches from location in that respect; the particular history remains
irreplaceable, and on the stronger claim its destruction remains a loss no
relocation of goods makes good, so what abundance removes is urgency and
location-binding, not the object. Miracleism in that world would supply participatory
significance of a different kind, incorporation without scarcity, with less
urgency rather than less meaning. \emph{A hostile inheritance.} Suppose a rare, intricately
coupled history perpetuates severe suffering, so that its continuation would be
the continuation of harm. The features that support care, the goods and the
relational reason, are then partly absent or outweighed, and the scarcity of
the structure does not rescue it: the bridge's reason concerns a continuation
carrying goods worth sustaining, and where the continuation bears mostly harm
the reason is defeated or outweighed. Bare persistence is not the bridge's object, and continued
existence is not sufficient for the reason; the continuation must carry what
made the history worth continuing. \emph{A final generation.} Suppose a rich
history will inevitably end and nothing can preserve it. The prospective
reason lapses with the agency, and with it the responsibility that depended on
leverage; the realized significance does not lapse, since nothing in the past
is undone by the end, and what remains fitting is acknowledgment, grief, and
the completion of what can still be completed. The Golden fruition is then
impossible, and the inheritance is in vain in the restricted sense
Section~\ref{sec:golden} defines, without having been without significance.

\subsection*{Defeaters}

Because the bridge is defeasible, a countervailing case does not refute it,
but the reader is owed an account of what defeats or limits its reason. The
antecedent can fail: an agent who has not adequately apprehended the object,
who has taken in the probability and not the structure, does not recognize
the reason and cannot act on it in deliberation; the reason itself, being a
fact about the relation, does not depend on being recognized, so what the
antecedent governs is access to the reason and not its existence. The reason
can be outweighed: present suffering, justice among the living, and
obligations to particular persons are reasons of their own and are not ranked
below continuation by anything in the framework; the bridge supplies a reason
and not a ranking, and Section~\ref{sec:selfaware} gives a case in which the
ranking is the whole difficulty. The reason can be defeated by quality: a
continuation that would perpetuate severe harm without prospect of remedy
loses the goods that gave the reason its content. The reason can lapse with
agency: where nothing can be done, nothing is required, and acknowledgment
takes responsibility's place. And the reason never licenses a means by itself:
that the continuation matters settles neither which actions preserve it nor
who should bear their costs, which is why the bridge is a premise about reasons
and not a decision procedure.

\section{The end-Cretaceous day: a tragic inheritance}\label{sec:kpg}

The exhibit is where affirmation and loss must be held together. There is one
day in the record that is, on the Nexic companion's
accounting, among the largest concentrated nodes of compounded improbability
in the route to us \cite{Nexus}, and it is the day the exhibit is built on.

The ordinary description first, and nothing in it is rejected. An impactor
struck a shallow, sulfur-rich sea, and the target's chemistry contributed to
the atmospheric forcing associated with the mass extinction. The dominant
dinosaurian regime ended: about a hundred and thirty-five million years of
ecological dominance measured from the end-Triassic, out of a reign of roughly
a hundred and sixty million years measured from the group's origin. Small,
fragile, endothermic mammalian lineages, among them the ones that would carry
the primate line, survived in enough number and in enough places to radiate
into the vacated world. Every clause is natural history with an ordinary
causal explanation.

Now the same day as a constraint problem. Terminate a globally entrenched
ecological regime that had excluded large-bodied mammals from every dominant
niche for the length of the Mesozoic, without terminating the small mammalian
channel that runs through it; do it in one event; do it with a body whose
angle, speed, and target chemistry are jointly such that the aerosol loading
is severe enough for the first and not so severe as to foreclose the second.
Human intelligence, given the stakes and the tools of the present, could
scarcely engineer the realized solution on purpose. That is what the phrase
\emph{computational divinity in the rocks} means: a performance metaphor, not
an agent claim, and not a complexity-theoretic one, since no input size,
search space, or objective has been defined. Natural history instantiated a solution whose
deliberate design would be staggeringly difficult, and the solution is in the
rocks, where it can be excavated, measured, and measured again. Nothing was
revealed. It was dug up.

The evidence for the pieces is local, and the ledger keeps it so.

\begin{center}\small
\begin{tabular}{@{}p{0.47\linewidth}lp{0.36\linewidth}@{}}
\toprule
Claim & Source & Scope\\
\midrule
Roughly thirteen percent of Earth's surface held hydrocarbon-rich sedimentary targets which, on the study's modeled soot mechanism, would have produced extinction-scale climatic forcing from a Chicxulub-scale impact & \cite{Kaiho2017} & the impact-location component on one modeled mechanism only; trajectory, velocity, angle, ecological selectivity, and the later evolutionary route are separate constraints\\
Impact-released sulfur of $67\pm39$ gigatonnes, several-fold below a prominent estimate, favoring a milder sulfur-driven winter & \cite{Nature2025Sulfur} & revises one component of the impact-winter mechanism; the exhibit survives changes in the aerosol decomposition if the coupled structure survives\\
Lower extinction rates among burrowing mammalian lineages under harsh, fluctuating climates & \cite{Burrow25} & one survival trait, independently supported\\
Impact in boreal spring; underground sheltering and dormancy discussed as possible survival mechanisms in the immediate aftermath & \cite{During22} & timing established; the refugia proposed, not established\\
The joint probability of the whole route is no larger than any single factor's and, under the companion's assumptions on how the factors couple, far smaller; the trait list read as an integrated preparation & \cite{Nexus} & the Nexic companion's structured reading; the strict reduction is its claim, since strongly coupled conditions can share a probability\\
\bottomrule
\end{tabular}
\end{center}

\noindent No entry establishes the probability or optimality of the entire
route; the joint claim is the companion's, conditional on its coupling
assumptions, and it is the one the exhibit interprets. On the Dark
Era reading of the Nexic companion \cite{Nexus}, the Mesozoic mammalian
condition, small, subordinate, nocturnal, denning and burrowing prey at the
margins through a hundred million years of disadvantage, is read as an
integrated preparation whose signs reverse under the day's realized lethality
profile, a continental-scale thermal pulse followed by impact winter: small
body size reduces energetic requirements and increases refuge flexibility,
marginal and low-trophic occupancy reduces dependence on the collapsing
dominant food web, nocturnality may carry traits useful under prolonged low
light, and burrowing shields from the pulse. The traits that marked the
mammals' inferiority under the ordinary regime are the survivors' kit under
the realized event, and the incumbents' formidability is what could not fit
down the burrow. Exaptation is the right comparator for traits acquiring new uses,
and it does not by itself settle the stronger interpretation; what the Nexic
reading adds is the claim that the trait list reverses sign together, as a
conjunction, under exactly the profile the conditioned law overweights.

The third description is the companions'. Let $\PP$ be the
unconditioned law over histories, $Q_T=\PP(\,\cdot\mid G_T)$ the law
conditioned on survival to horizon $T$, and $Q=\lim_TQ_T$ the limiting
conditioned law where the survivorship companion has established it. For any
event $E$ of the record and any finite horizon $T$,
\[
\frac{\PP(E\mid G_T,C)}{\PP(E\mid C)}=\frac{\PP(G_T\mid E,C)}{\PP(G_T\mid C)},
\]
so that if the end-Cretaceous configuration is enormously important to getting
from the Mesozoic to the Golden continuation, the whole joint configuration of
the day is enormously overweighted under the target-conditioned law. That much
is selection, available under the weaker descriptive claim. The Route World
hypothesis adds the stronger reading: if the law governing the realized
ensemble is the conditioned law, the day is not an event that later turned
out to be useful but a local realization of the conditioned law itself, what
conditioning looks like when instantiated in natural history. Conditioning can
instantiate coordination without a coordinator, and the exhibit is what that
looks like from inside the record; no mind aimed anything, and no backward
efficient cause acted.

The philosophical question the exhibit raises is the one the thesis cannot
avoid. Can we affirm our existence while refusing to treat the suffering and
destruction involved in its history as thereby justified? The day was, for the
incumbents and for most of the living world, a catastrophe without remedy: a
thermal pulse, a dark winter, and the extinction of some three quarters of
species, among them lineages of a hundred and sixty million years' standing.
Two words the exhibit invites, \emph{liberation} and \emph{surgical strike},
pull the reader toward the eventual beneficiary's standpoint, and even
without an intentional coordinator they can suggest that the losses were
justified by what followed. The first is kept and scoped: liberation names
the mammalian channel's release from a regime that had excluded it, read from
the channel's standpoint and from no other, and it describes one line through
the day and not the day. The second is dropped, since it implies aim.

Four attitudes come apart here, as they did in Section~\ref{sec:sublime}.
Gratitude for being alive. Recognition of historical dependence, that the
present lineages exist because that day took the form it did. Admiration for
the structure of the process, the form of the solution to the constraint
problem. Approval of everything the process involved. The first three are
fitting, on the paper's premises; the fourth does not follow from them and is
withheld. Being a beneficiary confers neither retrospective vindication of
every cause nor permission to disregard what was lost. The inheritance is
valuable and tragic at once, and nothing in the framework requires the tragedy
to be discounted so that the value can be affirmed. This is the register in
which the framework's reading of the day should be heard: the conditioned law
overweights the configuration and does not make it good, and a reader who
holds the beneficiary's gratitude and the mourner's grief together has
apprehended the object more completely than one who holds either alone.

Gratitude-like recognition without a giver can now be given content. What the
attitude acknowledges is dependence on a benefit that was neither earned nor
given: the present position exists because of a passage no one arranged and
that cost what it cost. It places its holder in no debt, since there is no one
to be indebted to, and it entails no obligation to a benefactor. What it does
is make the continuation of the benefit salient in the way gratitude
ordinarily does, by presenting the benefit as something received rather than
something owed, so that squandering it registers as a failure of
acknowledgment rather than as a neutral choice among options. That is the
only route by which, on this account, recognition motivates care.

\section{Futurity and finitude: the Golden Point}\label{sec:golden}

The Golden companion enters late, because the existential impact is already
present before it does. One looks backward at the record and asks what on
Earth this thing is; that question does not need the Golden Point. What the
Golden Point changes is one's relation to the thing: one is holding part of
its unfinished edge.

The Golden Point is the framework's specified viability condition and nothing
looser: the state in which the biosphere's external hazard is driven to zero
and the accumulated improbability of the route survives into a regime of
positive viability, defined through the survivorship model and its horizon
family \cite{GoldenSurvivorship,GoldenRouteWorld}; comparisons with
existential-risk discourse or with cultural visions of a desirable future do
not establish equivalence (Section~\ref{sec:selfaware}). The division of
labor is stated as the Miracleist interpretation and not as a result. On that interpretation,
the Miracles source the inherited significance; Golden continuation permits
its living fruition; and Departure terminates the inheritance before that
fruition. What Departure cannot do is unmake the realized structure: the
end-Cretaceous day does not cease to have happened, or to have the structure
found significant, if the lineage ends tomorrow. The past gives the Miracles,
whose coupled closure is the inherited structure; the present is a section of
the same continuation whose previous sections crossed the bottlenecks; the
future divides into continuation and termination. The Golden Point does not
retrospectively create the significance of the end-Cretaceous passage or of
abiogenesis; it controls whether the route reaches durable fruition. The word \emph{Miracle} carries the companion's
technical sense throughout, a target-amplified, independently specified,
coupled configuration, and no religious sense \cite{GoldenRouteWorld}.

Scheffler is the nearest interlocutor, and the comparison should be made at
his resolution rather than as a broad opposition between his future and this
paper's past. His collective-afterlife argument distinguishes ways in which
present valuing depends on the existence of future people, attitudinal,
evaluative, and justificatory, and argues that much of what matters to us now
would lose its point, and in some cases its value, if humanity were to end
soon after our deaths \cite{Scheffler13}. Miracleism agrees about the
direction of the dependence and differs about its object. The dependence it
accepts is evaluative and concerns the continuation: whether the inheritance
reaches fruition depends on the future, and the value of present stewardship
depends on that possibility. The dependence it denies concerns the source: the
realized structure's significance does not depend on the future, because it is
already realized. The contribution is therefore an account of inherited value
and prospective fulfillment, two objects with different temporal homes, which
Scheffler's single notion of the collective afterlife does not need to
separate and this paper must.

The separation forces a question the phrase \emph{in vain} can obscure. If
Departure does not erase the history's significance, what exactly has been in
vain? In vain names the failure of the inherited structure to
reach living fruition, never the retroactive erasure of its realized
significance; but why does that fruition matter, and which achievements
remain valuable without it? Consider partial achievement. A flourishing but
finite civilization that realized understanding and worthwhile experience
would have realized goods of the inheritance, and those goods, once realized,
are not undone by a later end; a reader who takes the paper to say that a
finite history is worthless has misread it. What Golden continuation adds is
not more value of the same kind. It removes the condition that otherwise
makes every realized good's continuation a matter of luck: before the Golden
Point the continuation of the goods remains exposed to the external hazard the
route has so far survived by passages no one controlled, and after it the
continuation is secured against that hazard, the inheritance held rather than
wagered. Fruition in the paper's sense is that conversion, from a continuation
that persists by hazard to one that persists by viability, and it matters
because it is the only form in which the inheritance can be kept rather than
enjoyed while it lasts.

What continues is also to be stated rather than settled by a metaphor. The
Golden target as the companions specify it is the biosphere's continuation,
the external hazard driven to zero for the system as a whole; the Miracleist
object of care is the integrated history whose current carrier is that
biosphere, including its observer-bearing lineages and the understanding they
have accumulated. These overlap without being interchangeable. A biosphere that persisted without
observers would preserve the carrier and lose the incorporation relation;
accumulated understanding cannot persist without carriers; a particular
lineage is not the structure, since the structure has already passed through
the replacement of lineages. Which of these sub-continuations would count as
fruition and which as partial loss is a question the framework's own object
poses and this paper does not answer; it records that the continuation the
care bridge concerns is the goods-bearing one, so that a continuation stripped
of the goods would not be the object the reason was about.

Strong Route World is where the stronger ontology contributes without being
reduced to probability bookkeeping. The hypothesis is that the governing measure itself is
conditioned in the relevant way, so that the target is a physically privileged
condition, the one under which the realized record is typical rather than
freakish. A physically privileged target is not thereby a goal agents have
reason to adopt, any more than a physical equilibrium is good because systems
tend to it. What Strong Route
World contributes to the reasons is threefold. It identifies which
continuation is at issue, the one the record's own law privileges, so that
care has a determinate object. It makes the relation between the inherited
structure and that continuation non-accidental: under the hypothesis the
structure is a realization of the conditioned law rather than a lucky run
that happens to point somewhere, which licenses reading inheritance and
continuation as one object with a past and a future rather than as a past and
an unrelated hope. And it makes the continuation prospectively testable, which
turns stewardship into participation in an inquiry rather than devotion to a
creed. Whether that
privileged condition gives reasons still requires the care bridge; the
hypothesis changes the object's modal status, and the bridge says why an
object of that status, with those goods and that relation to us, matters.
Withhold the hypothesis and the weaker interpretation remains, which is most
of the paper; accept it and the continuation becomes the condition the
inheritance's own law marks.

Nagel's challenge returns here in temporal form, and the answer is the same:
a longer future or a larger encompassing project does not justify present
concern by its size, and what justifies it, if anything does, is the
particular value carried and the particular relation, the goods ours by
constitution, their continuation without substitute, and ourselves placed
where it can be lost. A generation with material leverage over entry into the
positive-viability region inherits neither cosmic rank nor a guarantee of
success, but responsibility at an unusually high-leverage point. The
civilization-and-nature companion gives one exact schedule model in which the
retained continuation option is nondecreasing as gates are cleared
\cite{GoldenCivNature}, which is not a law that every later moment carries
greater stakes than every earlier one; the narrower point is that where
cleared prerequisites change the live option set, forfeiting the route
forfeits an option that did not exist before those gates were crossed. That is the winner's irony in its
properly scoped form: where a proposal would forfeit the continuation option
under the model, wishing the route dismantled for the living world's sake
pushes away the continuation through which the inheritance remains live,
whatever the wisher's reasons. That a given proposal does so is a result to be
established for that proposal, not a verdict on its advocate. Individual mortality is untouched
throughout; the fruition at issue is collective, and the moments remain mortal
exactly as the Route World companion describes them.

\section{Public implications, and their limits}\label{sec:selfaware}

The incorporation relation has a collective form, where the discovery, if
secured, stops being something that happens to a reader and becomes something
that could happen to a species. Three topics belong here: what it would be for
a species to possess the concept, how responsibility attaches without a
certificate of worthiness, and how the account could enter public life.

\subsection*{Golden species self-location}

The underlying property is the Nexic companion's: an entity is self-aware when
it can represent its own situation as its own, including what it would be for
that situation to be won or lost \cite{Nexus}. The companion's discipline
transfers one scale up: behavioral
proxies for the property are category-error shaped, so no species-scale proxy
is proposed here either. Lifting the property to the species gives the
criterion.

\begin{definition}[Golden species self-location; \textbf{[I]}]\label{def:selfaware}
A species has \emph{Golden species self-location} when the concept of its
own Golden and Departure poles, its modeled winning and losing conditions
represented as conditions of its own continuation, is available to it for
shared public reasoning: held not merely by some of its members but in a form
that can enter collective deliberation about its continuation.
\end{definition}

\noindent The modifier \emph{Golden} marks a narrow representational
achievement, collective self-location between the framework's own modal
poles, and not a claim about self-awareness in the ordinary philosophy-of-mind
sense. Three levels come apart: representation by individual members,
availability for shared public reasoning, and incorporation into effective
institutional deliberation. The first is necessary and is not the criterion,
since this paper represents the poles and its existence does not satisfy the
definition; the criterion is the second; the third is a further achievement
the second does not guarantee. By the criterion so stated, humanity has
not passed: it acts at scale with planetary consequences, it has collective
understandings of ecological and existential stakes, and these particular
modal poles are not yet established objects of its shared deliberation.

Two distinctions bound the claim. Aggregate planetary effect is not
organized collective agency: consequences at the scale of the biosphere do
not make the species an agent with a unified representation and decision
procedure, which in the demanding sense requires an organization that
aggregates judgments and acts on them \cite{ListPettit11}, and nothing of
that kind is claimed; the criterion concerns a representation shared among
the many agents that exist. And a newly named concept must earn its novelty
by content. Climate, nuclear-risk, and existential-risk discourse already
represent collective hazards and losing conditions. What Golden self-location
adds is an object and a relation: a specified viability condition rather than
hazard reduction without a terminus, and the location of the present within
an inherited structure whose continuation that condition concerns. A species
that reduces hazards without representing what the reduction is for has
executive competence without self-location; one that represents the poles
can own its conduct \cite{GoldenCivNature}, and under the Route World
hypothesis target-consistent conduct before target-representation is what a
Golden-conditioned record should show \cite{GoldenRouteWorld}. The absence of the concept has a recognizable
signature, \emph{value completion before warrant}: first-order commitments
asserted as worthy with no external object supplied that adjudicates their
ordering when they conflict. Possession of the concept is not endorsement of
the hypothesis; a species can represent its poles and remain uncertain whether
its record is drawn under conditioning on one of them, and that uncertainty,
held correctly, is itself a form of the awareness the criterion names. Shared
continuation stakes establish neither a single authorized decision-maker nor
equal responsibility, which the next subsection takes up.

\subsection*{Responsibility without metaphysical worthiness}

The argument does not require a verdict that humanity, a generation, or any
person is metaphysically worthy of the position occupied, and Miracleism
rejects that variable as a primitive. What is rejected is specific: a
person-level moral essence, a hidden state of worthiness stamped onto a
person independently of what that person does and consulted before
responsibility can attach. Ordinary character ethics is not committed to it;
evaluating persons by stable dispositions and histories is legitimate, and
what is rejected is the further step that turns the summary into a
substance.

Rejecting the primitive settles less than it might seem to, and five things it
leaves standing should be distinguished: basic moral standing, the status in
virtue of which a being's interests count, which is not earned; character,
real and assessable as a pattern; desert, which attaches to conduct and its
consequences; accountability, which attaches to agents for what they have
done; and capacity, what an agent can affect, which varies by person and
circumstance. The moral grammar the paper proposes begins at the level of
choices and trajectories, better and worse uses of an available option, with
person-level labels as compressed descriptions rather than certificates. A
person can make a terrible choice and be answerable for it without becoming a
bad kind of thing; repeated choices entrench character without any essence
settling the next choice in advance. At the collective scale the same point blocks two symmetrical
fatalisms: ``humanity is good'' supplies no guarantee of continuation, and
``humanity is bad'' supplies no excuse for Departure.

Responsibility, on this account, follows from the combination of a relevant
relationship and situated agency, and its extent depends on what an agent can
reasonably do. The relationship is the one Section~\ref{sec:sublime} defended,
being a continuation of the inheritance; the agency is the leverage an agent
actually has over the continuation's course, which differs enormously across
persons. Someone with little capacity to influence collective outcomes does
not thereby occupy a less significant life: the significance is participatory
and shared, the responsibility differentiated by capacity. That is why the
framework replaces the question whether we are worthy of this with the
question what, given that we are here, the good choice from here is, a
relocation of responsibility rather than its relaxation.

The relocation does not identify a party, leader, institution, or policy as
Golden or anti-Golden, and a concrete conflict shows why. Suppose a proposal
claims to improve the biosphere's distant viability, by a large program of
hazard reduction, say, while imposing serious present costs on particular
people who did not choose them. Three steps are distinct. The Golden model may
establish something about the means, whether the program bears on the
viability condition at all and by how much, an empirical question answered by
ordinary evidence. The care bridge contributes a reason to take the
continuation seriously and nothing about the distribution of costs. Licensing
the sacrifice requires further premises, about fair allocation, consent, the
claims of the present against the future, and the authority of whatever
institution would impose it, which the framework does not supply and which
reasonable people who share the object will dispute. Agreement on the object
requires no uniformity in the judgments that follow: Miracleism supplies the
stakes and the object of concern, not a factional authority.

\subsection*{Public expression}

How the account could enter public life is a question about culture, and it is
treated here as a bounded extension with its own grade: the claims are
hypotheses about reception, not consequences of the argument above.

Four conditions that the vocabulary of vacancy runs together should be
separated, since they occur independently and require different evidence:
inadequate warrant for a shared frame of significance; unmet personal demand
for one; absence of a shared public understanding where individuals possess
one; and the inability of institutions to act on an understanding they
possess. The claim is not that secular people lack meaning. The narrower
hypothesis, the \emph{Cultural Vacancy Hypothesis}, is that particular
publics may find that their existing frameworks do not
adequately supply one kind of orientation: an external account that locates
human life at once within natural history, cosmic scarcity, and collective
futurity, in a structure not authored to satisfy the need it answers.
Prevalence is a question for research, not a premise. And the niche is not
empty. Religious naturalism already offers reverence toward the natural world
grounded in the scientific account of it, Goodenough's being the fullest
expression \cite{Goodenough98}; Big History offers a shared narrative from the
origin of the universe to the present as a frame for collective
self-understanding \cite{Christian04}; environmental and existential-risk
movements supply practical orientation toward continuation. What Miracleism
would add to each is the specified object, the incorporation relation, and
the viability condition, and whether that addition answers a demand the
others leave unmet is what reception would test.

Reception itself is not one thing. Visibility, comprehension, endorsement,
motivation, and durable practice are distinct achievements; digital
circulation can raise the first without improving the second, and widespread
familiarity need not produce agreement or action. A cultural resource changes
conduct, on Swidler's account, by entering the repertoire from which people
construct strategies of action, which belief alone does not accomplish
\cite{Swidler86}. Forecasts of uptake belong to a separate research
prospectus; what is retained is the possibility of shared public
understanding, without a forecast of its extent.

One hypothetical pathway makes the possibility concrete without predicting
it. An organization responsible for long-horizon planning learns the
conditional account, asks whether the care bridge's reason, if accepted,
bears on its remit, and examines whether that reason justifies revising an
existing practice, say the horizon over which it discounts low-probability,
high-consequence hazards. An actual change is a documented alteration of
practice traceable to that evaluation, and a sequence of such evaluations
that leave practice unchanged counts against the hypothesis. Jasanoff and
Kim's sociotechnical imaginary, a collectively held vision of a desirable
future attainable through science and technology, names what a public
realization of Golden reasoning would be \cite{JasanoffKim15}; it describes
the expression and leaves the account of its object intact. The philosophical
argument of this paper is assessable whether Miracleism becomes a specialist
interpretation, inspires varied practices, or gains no attention at all.

\section{What Miracleism claims, and what would undermine it}\label{sec:ladder}

The claims above have different statuses and different defeaters, and they
are collected here so that the paper is not read as deriving significance
from Bayes' theorem. It does not. It states what the world would be like,
existentially, if independently testable claims were true and two
philosophical premises were accepted. The core interpretation rests on the
awe and care bridges; the stronger thesis additionally adopts the
organic-unity premise of Section~\ref{sec:sublime}, and the table lists it
separately so that its rejection is seen to cost exactly that thesis.

\begin{center}\small
\begin{tabular}{@{}lc@{}}
\toprule
Rung & Grade\\
\midrule
Accordance, calibrated typicality under the draw & [P]\\
Rare-Earth inference under stated antecedents & [P/C$_{\mathrm{emp}}$]\\
High-Contrast Earth and Optimal Earth & [C]\\
Exotic-structure reading of the coupled closure & [I]\\
Strong Route World: the record under the conditioned law & [C]\\
Fitting-awe bridge (Premise~\ref{prem:awe}) & [$\Phi$]\\
Care bridge (Premise~\ref{prem:care}) & [$\Phi$]\\
Organic-unity premise, restricted to unities of goods (stronger thesis only) & [$\Phi$]\\
Miracleism as existential interpretation & [I]\\
Golden species self-location as a public possession & [I]\\
Cultural Vacancy Hypothesis and public uptake & [H]\\
\bottomrule
\end{tabular}
\end{center}
\noindent Here [P] is proved, [C] conjectural or conditional, [I]
interpretive, [$\Phi$] a philosophical premise defended but not derived, and
[H] an empirical hypothesis about reception. The claims have different
possible defeaters, and they are listed. If the Rare-Earth inference fails
badly, the warrant for the scarcity component is withdrawn; that withdrawal
does not establish an occupied cosmos, and what remains is an inheritance
whose scarcity is unknown, which weakens the location-bound reason of
Section~\ref{sec:sublime} toward the abundant-world case without touching
incorporation. If the purported Miracles dissolve under better priors, or
their guards are shown violated, rarity fails. If the couplings disappear
under preregistered tests, exotic structure fails and what remains is a list.
If High-Contrast Earth or Optimal Earth fails, the concentration reading fails
or weakens, while the end-Cretaceous day may remain rare and target-enriched.
If Strong Route World fails, the stronger ontic claim fails: the event may
remain selected without being conditioning instantiated in the record, and the
interpretation falls back to the weaker layer of Section~\ref{sec:golden}. The
bridges can fail with every empirical claim intact, which is the point of
listing them. The fitting-awe bridge fails if awe's presentation of its object
is shown not to match the object described, or if fittingness requires more
than accuracy. The care bridge fails if relational reasons from unchosen
constitution are rejected in general, or if the goods route is found
insufficient for location-bound care. The organic-unity premise fails if non-additive value is
rejected in general, or if its restriction to unities of goods is shown not
to hold of this history; its rejection defeats the stronger thesis and leaves
the core interpretation standing. The reception hypothesis fails if publics that comprehend
the account do not find that it supplies the orientation claimed, or if
comprehension produces no change in practice under evaluations of the kind
Section~\ref{sec:selfaware} describes; its failure leaves the philosophical
argument untouched.

One further fence guards the ladder from both sides. A deliverance of
perception, reason, memory, or moral judgment supplies prima facie support
that a theory may defeat, but the defeater has to contain information, naming
the premise, rule, domain condition, or inference that fails; an unspecified
deeper layer whose sole function is to make a conflict cease to count is
logically idle, the point being Popper's against ad hoc immunization
\cite{Popper63} and phenomenal conservatism's against undefeated-appearance
denial \cite{Huemer07}. The valley invites the error in both directions: one
can protect an inherited source of meaning from contrary evidence by inventing
a safety layer, and one can equally protect nihilism by declaring in advance
that no natural structure, however independently discovered, could count as
meaning proper. Miracleism permits neither move; the reversal has to be
purchased by a statable evidential object and explicit bridges.

Three literatures are neighbors: naturalistic objectivism, on which features
of the physical world can confer meaning independently of anyone choosing
them \cite{MetzSEP}; the scientific sublime \cite{Gross18}; and cosmic
purpose without the traditional God \cite{Mulgan15,Goff23}. Miracleism's
originality rests neither on science having previously produced only
disenchantment nor on naturalistic wonder being new. It rests on the proposed
object, its structure, and our relation to it: the process historically
responsible for disenchantment reveals, if the Nexic and Golden hypotheses
hold, an independently testable structure that reverses the evidential sign
without reversing the method and includes the reader among its continuation
states. The reversal, if it is one, must be earned by the standards that
produced the descent, and no tradition that lowers those standards is
Miracleism, whatever it calls itself.

\section{Conclusion}

If the empirical and conjectural rungs survive, then on the Miracleist reading
Earth is one of extraordinarily few observer-bearing worlds within the volume
that could ever be surveyed, and the route that brought it from its first
replicators to the reader is a knot rather than a sequence of accidents: a
coupled structure in which the end-Cretaceous day and the bottlenecks before
and after it stand where they stand because the others do, at a surprisal no
one can render and amplified by a target no one chose. That structure was
neither received by revelation nor authored by its discoverer. It was dug up,
by the same inquiry that removed the objects previously offered as meaning's
warrant. And the person who digs it up finds, last, that they are made of its
continuation: the vast thing contemplated contains the one contemplating it,
which is why what occurred here is too vast to be conveyed and too personal to
be named.

The paper's own claims are two bridges and their consequences, with one
further premise, that an organization of goods bears value beyond the goods,
carrying the stronger thesis. Awe is fitting
toward such an object because the object has what awe presents. Its
continuation gives its inheritors a defeasible reason for care, because the
goods it carries are irreplaceable within the surveyable volume, because the
inheritors are constituted by it, and because its continuation can still be
lost; the reason is outweighable, defeated by a continuation that would carry
only harm, and silent about means. The inheritance is valuable and tragic at
once, and affirming it requires no approval of what it cost. The Golden Point
adds the one thing that makes the finding a task rather than a museum: the
structure is not behind us. We are sitting inside its unfinished edge, where
the inheritance can still be held rather than wagered, and where
responsibility attaches to whoever can affect that, in proportion to what they
can do, without any prior verdict that they deserve to be there. Whether the
Miracles were in vain, in the one sense the word carries here, is not yet
decided. Miracleism is the name for holding that without possessing it.

\appendix
\section{The exotic-structure witnesses}\label{app:witnesses}

The witnesses of Section~\ref{sec:exotic} are defined here with the audits
that limit what they claim. \begin{definition}[Exotic structure, with higher-order witnesses; \textbf{[I]} on the
companions' \textbf{[C]}]\label{def:exotic}
Fix an admissible attribute basis $\mathfrak A_C$, chosen independently
of any coupling claim, and quotient it by the framework's validity
equivalence, $[A]=\{B:B\equiv^{V}_{C}A\}$, so that the vertices below
lie in $\mathfrak A_C/{\equiv^{V}_{C}}$. Fix a target attribute $A_j$
and a nonempty context set $S\subseteq V\setminus\{A_j\}$. For each
$i\in S$ write $A_i^{(1)}=A_i$ and $A_i^{(0)}=A_i'$, and for an
assignment $\epsilon_S\in\{0,1\}^{S}$ let
$\operatorname{Set}_{\epsilon_S}$ jointly settle every member of $S$
to the state selected by $\epsilon_S$. Define the target effect of
$A_j$ in that settled context by
\[
\Delta_j(\epsilon_S)=
\log\frac{\PP\bigl(G_T\mid
\operatorname{Set}_{\epsilon_S}(A_j,C)\bigr)}
{\PP\bigl(G_T\mid
\operatorname{Set}_{\epsilon_S}(A_j',C)\bigr)},
\]
computed under the framework's settling operator $\operatorname{Set}$
\cite{Aethus}; conditioning is not substituted for settling, for the
reason given in Section~\ref{sec:rarity}. The pure interaction of the
context $S$ on $A_j$ is the discrete mixed difference
\[
\iota_{S\to j}
=
\sum_{\epsilon_S\in\{0,1\}^{S}}
(-1)^{|S|-\sum_{i\in S}\epsilon_i}\,
\Delta_j(\epsilon_S).
\]
For $S=\{A_i\}$ this reduces to
$\iota_{\{i\}\to j}=\Delta_j(A_i)-\Delta_j(A_i')$, so the pairwise
witness is recovered exactly as $w_{ij}=|\iota_{\{i\}\to j}|$. The
dependency object is therefore a weighted directed \emph{interaction
hypergraph} $\mathcal H_{\Ec}=(V,\mathcal H)$, carrying a hyperedge
$S\to A_j$ whenever $\iota_{S\to j}\neq0$, with weight $|\iota_{S\to
j}|$. A nonzero hyperedge says that the joint settled configuration of
$S$ changes $A_j$'s target role in a way not exhausted by lower-order
interaction terms.

The finite interaction witness is used only on the positive-survival
settling cube, where every target probability entering the mixed
difference is positive. Structural zero-survival cases, in which a
settling forces $\PP(G_T\mid\cdot)=0$, enter $\Psi_{\Gc}$ as forcing
relations rather than as infinite hyperedge weights, so that absolute
incompatibilities belong to the closure structure and no difference of
infinities arises. At order $S$, the null is \emph{target-separability},
$\iota_{S\to j}=0$; full target-separability requires this at every
nonempty admissible order, not merely for singletons.

Let $\Psi_{\Gc}$ be the monotone forcing operator that adds the
attributes whose target-relevance commitments the current set forces,
and define
\[
F_{\Gc,A_{\mathrm{root}}}(A)
=A_{\mathrm{root}}\cup A\cup\Psi_{\Gc}(A).
\]
The vertex set of the exotic structure generated by the root is
\[
V(\Ec)=\operatorname{lfp}\!\bigl(F_{\Gc,A_{\mathrm{root}}}\bigr),
\qquad
A^{(0)}=A_{\mathrm{root}},\quad
A^{(m+1)}=A^{(m)}\cup\Psi_{\Gc}\bigl(A^{(m)}\bigr),
\]
the least fixed point containing $A_{\mathrm{root}}$, which for a finite
closure is $V(\Ec)=A^{(D)}$ at the first stabilization step; the exotic
structure itself is the structured closure
$\Ec=\bigl(V(\Ec),\ \mathcal H_{\Ec},\ \Psi_{\Gc}{\restriction}V(\Ec)\bigr)$,
the vertex set together with its interaction hypergraph and its forcing
relations, and the inequality $\Ec\neq\{A_1,\dots,A_n\}$ of
Section~\ref{sec:exotic} marks that difference between a structured closure
and a bare inventory. Its witnesses are the higher-order coupling
\[
K(\Ec)=\sum_{(S\to j)\in\mathcal H}|\iota_{S\to j}|,
\]
the depth
$D(\Ec)=\min\{m:A^{(m+1)}=A^{(m)}\}$ for a finite closure, the
synchronous iteration depth of the stated presentation of $\Psi_{\Gc}$, and
the rarity of the realized conjunction the closure represents,
$R(\Ec)=-\log\PP\bigl(\bigcap_{A\in V(\Ec)}A\mid C\bigr)$.
\end{definition}

\noindent Three audits accompany the definition, each limiting what it
claims. First, the null is higher-order target-separability and not
statistical independence, and pairwise tests are insufficient. Let $A,B,C$ be
independent fair bits and let a binary target $G$ be generated by a noisy
parity rule: for some $0<\epsilon<\tfrac12$,
\[
\PP(G=1\mid A,B,C)=
\begin{cases}
1-\epsilon,&A\oplus B\oplus C=1,\\
\epsilon,&A\oplus B\oplus C=0.
\end{cases}
\]
Marginalizing the un-set bit makes every singleton context pairwise
target-separable, while fixing the other two bits makes their joint mixed
difference on the third nonzero: a higher-order target coupling can be
invisible to every pairwise witness with all probabilities strictly positive,
and the hyperedge witness exists so that such synergy is not misclassified as
an edgeless list. What a nonzero hyperedge establishes is nonseparability of
the specified kind, the target role of $A_j$ not being recoverable from the
members of $S$ taken one at a time. The word \emph{non-reductive} below
carries exactly that content; no emergence beyond nonseparability, no downward
causation, and no irreducibility of the members' own physics is asserted.

Second, depth is a property of the stated presentation and not an intrinsic
height. The synchronous iteration counts rounds of parallel derivation, which
follow the earliest available derivations rather than the longest: with the
irredundant forcings $r\Rightarrow a$, $a\Rightarrow d$, $r\Rightarrow b$,
$b\Rightarrow c$, $c\Rightarrow d$ from the root $r$, the iteration gives
$A^{(0)}=\{r\}$, $A^{(1)}=\{r,a,b\}$, $A^{(2)}=\{r,a,b,c,d\}$, so $D=2$,
while the chain $r\Rightarrow b\Rightarrow c\Rightarrow d$ has length three.
Reachability under an acyclic forcing relation is a strict partial order; the
immediate forcing relation itself need not be transitive, and $D$ is computed
on it. The iteration depth coincides with the longest chain in some
structures, for example when the acyclic forcing structure is a rooted tree
with one supporting predecessor per non-root attribute, and the coincidence
is not confined to them: the diamond $r\Rightarrow a\Rightarrow d$,
$r\Rightarrow b\Rightarrow d$ has depth two and longest chain two. In
general $D$ is reported with the presentation it is computed on. Nothing in
the paper's argument turns on its numerical value; what is used is that
$D>0$, that the closure is not reached at the root.

Third, exotic grade is defined only on description classes modulo
refinements that preserve target information, and the quotient is constructed
for one case and stated as a requirement for the rest.
If a basis attribute $A$ is refined into $A_1,\dots,A_m$ and the refinement
carries jointly no target-relevance information beyond $A$, the two
descriptions are identified by a merge: refinement cells whose settled
configurations induce the same target effects $\Delta$ in every admissible
context are identified, and the merged attribute's event is the union of the
identified cells, with probability their sum. The aggregation rule is what
makes the rarity descend: refine $A$ by an independent nuisance attribute $N$
carrying no target information, so that the cells $A\cap N$ and
$A\cap N^{c}$ have the parent's effect profile; each cell alone has a larger
surprisal than $A$, but the merged event is $(A\cap N)\cup(A\cap N^{c})=A$,
its probability $\PP(A\cap N\mid C)+\PP(A\cap N^{c}\mid C)=\PP(A\mid C)$,
and the parent's surprisal is recovered. $K$ and $D$ are then computed on the
merged basis, on which the identified cells are one vertex with one effect
profile, so that duplicated interaction mass is aggregated rather than
multiplied by redescription. What this establishes is invariance under
refinement by target-irrelevant attributes; invariance of all three witnesses
under every admissible redescription is a requirement on admissible
descriptions whose full construction belongs to the companions, and a basis
on which the witnesses fail to descend is inadmissible. The raw hyperedge sum
on an unmerged basis establishes nothing about invariance; this is the
granularity guard the microdescription guard on rarity requires here too.

\section*{Contribution disclosure}
\addcontentsline{toc}{section}{Contribution disclosure}

\noindent The framework, its premises, and the thesis are the author's. The
philosophical development of Sections~\ref{sec:construction}--\ref{sec:selfaware}
was drafted with the assistance of Claude Fable~5.1 to the author's brief,
from the author's premises and the referee assessments it answers; the
diffusion hypotheses and the dated timing conjecture of earlier drafts are
held in a separate research prospectus. Responsibility for every claim is
the author's.

\printbibliography
\end{document}